% GENERATED FILE. Edit canonical lessons, then run npm run generate:books.
% canonical-source-hash: fnv1a64:e720a79a4f61a164
% canonical-lessons: ZH-C326-guojia, ZH-C326-shijie, ZH-W326-jie-boundary, ZH-C326-waijie, ZH-C326-jianli, ZH-C326-wenming

\chapter{Around the World}
\label{ch:zh-around-the-world}
\hlchaptermodality{326}

\begin{chapteropening}
\textbf{By the end of this chapter:} \emph{I can talk about countries, the world, the outside world, founding something and old civilizations.}

Complete the last lesson of chapter 326: I can talk about countries, the world, the outside world, founding something and old civilizations.
\end{chapteropening}

\section[\texorpdfstring{\zh{国家} (guójiā)}{国家 (guójiā)}]{\zh{国家} (guójiā) — country; nation}

\label{lesson:ZH-C326-guojia}



\begin{quote}
Before the new one: say the Mandarin for for generations on end, then the Mandarin for to miss someone dearly.
\end{quote}

\subsection*{You'll want to know: \zh{国家}}
\textbf{\zh{国家}} — \emph{guójiā} — \textquotedblleft{}country; nation\textquotedblright{}.

\begin{quote}
\textbf{\zh{我去过三个国家}} \emph{wǒ qù guo sān ge guójiā} — I've been to three countries.
\end{quote}

\begin{cousinweb}
\textbf{\zh{国家}} (\emph{guójiā}) is built from \zh{国} 'a state' and \zh{家} 'a home', both inherited from Middle Chinese, \emph{kwok} and \emph{kae} (Guangyun, in Baxter's transcription), which glosses the first 'a realm' and the second 'a dwelling'.
\end{cousinweb}

\begin{grammarlens}[title={What you've built}]
One more word for telling what happened.
\end{grammarlens}

\subsection*{Your turn}
\begin{itemize}
\raggedright
  \item \emph{Say it:} \emph{guójiā}
  \item \emph{Say it:} \emph{guójiā}, once more
  \item \emph{Say it:} \emph{wǒ qù guo sān ge guójiā} — I've been to three countries.
  \item \emph{Recall:} say the Mandarin for for generations on end, then the Mandarin for to miss someone dearly, then say \emph{guójiā} again
\end{itemize}

\begin{tcolorbox}[breakable,colback=teal!4,colframe=teal!35!black,title={Before you move on}]
What does \zh{国家} mean? (\textquotedblleft{}Country; nation\textquotedblright{}.) Say it once more.
\end{tcolorbox}
\section[\texorpdfstring{\zh{世界} (shìjiè)}{世界 (shìjiè)}]{\zh{世界} (shìjiè) — world}

\label{lesson:ZH-C326-shijie}



\begin{quote}
Before the new one: say the Mandarin for to miss someone dearly, then the Mandarin for a nation with its own government.
\end{quote}

\subsection*{You'll want to know: \zh{世界}}
\textbf{\zh{世界}} — \emph{shìjiè} — \textquotedblleft{}world\textquotedblright{}.

\begin{quote}
\textbf{\zh{我想去看世界}} \emph{wǒ xiǎng qù kàn shìjiè} — I want to go and see the world.
\end{quote}

The new character, in pieces: \zh{界} — nine strokes: on top a field, a box of five strokes cut into four; under it \zh{人} opened wide, and under that a short stroke falling to the left and a vertical. The next lesson puts it in your hand.

\begin{cousinweb}
\textbf{\zh{世界}} (\emph{shìjiè}) is built from \zh{世} 'a generation' and \emph{jiè} 'a boundary', both inherited from Middle Chinese, \emph{syejH} and \emph{keajH} (Guangyun, in Baxter's transcription), which glosses the first 'a generation' and the second 'a border' and 'an edge'. The sense 'the world' is not given there.
\end{cousinweb}

\begin{grammarlens}[title={What you've built}]
One more word for telling what happened.
\end{grammarlens}

\subsection*{Your turn}
\begin{itemize}
\raggedright
  \item \emph{Say it:} \emph{shìjiè}
  \item \emph{Say it:} \emph{shìjiè}, once more
  \item \emph{Say it:} \emph{wǒ xiǎng qù kàn shìjiè} — I want to go and see the world.
  \item \emph{Recall:} say the Mandarin for to miss someone dearly, then the Mandarin for a nation with its own government, then say \emph{shìjiè} again
\end{itemize}

\begin{tcolorbox}[breakable,colback=teal!4,colframe=teal!35!black,title={Before you move on}]
What does \zh{世界} mean? (\textquotedblleft{}World\textquotedblright{}.) Say it once more.
\end{tcolorbox}
\section[\texorpdfstring{\zh{界} (jiè)}{界 (jiè)}]{\zh{界} — a field over \zh{人}}

\label{lesson:ZH-W326-jie-boundary}



\begin{quote}
The word first, then the character inside it.

\begin{itemize}
\raggedright
  \item \emph{Recall:} say \emph{shìjiè} — \textbf{R1}, one lesson back, before it decays
\end{itemize}

\begin{quote}
\zh{人} \zh{世}
\end{quote}

\emph{Read it:} each one aloud, with nothing said first — \emph{rén}, \emph{shì}
\end{quote}

\subsection*{Script — a field over \zh{人}}
\hlblockfigure{\detokenize{figures/ZH-W326-jie-boundary-filmstrip.pdf}}{How \zh{界} is written, stroke by stroke}

\begin{quote}
\zh{界}
\end{quote}

\textbf{Two pieces, top to bottom.}

\textbf{The top piece, a field, five strokes.} A short vertical on the left. A bar that runs right and turns down. A bar across the middle, then a vertical down the middle. Then the bar that closes the foot.

\textbf{The bottom piece, four strokes.} First \zh{人}, opened wide: a long stroke falling to the left, then a long stroke falling to the right. Under them, a short stroke that falls to the left, and a vertical on the right.

5 + 4 = nine.

\subsection*{Your turn}
\begin{itemize}
\raggedright
  \item \emph{Write it:} \zh{界}, all nine, with the model in front of you
  \item \emph{Cover:} write it again from memory, and count the pieces, not the strokes
\end{itemize}

\begin{tcolorbox}[breakable,colback=teal!4,colframe=teal!35!black,title={Before you move on}]
How many strokes in \zh{界}? (\textbf{Nine}.) What is on top? (\textbf{A field, five strokes}.) What opens the bottom piece? (\textbf{\zh{人}}.)

Source: \href{https://raw.githubusercontent.com/chanind/hanzi-writer-data/68d10a4b21150cae5e1ebbd223eed289cf32d90c/data/\%E7\%95\%8C.json}{Hanzi Writer Data \zh{界}, pinned stroke paths and medians}.
\end{tcolorbox}
\section[\texorpdfstring{\zh{外界} (wàijiè)}{外界 (wàijiè)}]{\zh{外界} (wàijiè) — the outside world}

\label{lesson:ZH-C326-waijie}



\begin{quote}
Before the new one: say the Mandarin for a nation with its own government, then the Mandarin for the whole wide world.
\end{quote}

\subsection*{You'll want to know: \zh{外界}}
\textbf{\zh{外界}} — \emph{wàijiè} — \textquotedblleft{}the outside world\textquotedblright{}.

\begin{quote}
\textbf{\zh{外界都很关心他的病}} \emph{wàijiè dōu hěn guānxīn tā de bìng} — People everywhere are concerned about his illness.
\end{quote}

\begin{cousinweb}
\textbf{\zh{外界}} (\emph{wàijiè}) is built from \zh{外} 'outside' and \zh{界} 'a boundary', both inherited from Middle Chinese, \emph{ngwajH} and \emph{keajH} (Guangyun, in Baxter's transcription), which glosses the first 'outside' and 'far' and the second 'a border'.
\end{cousinweb}

\begin{grammarlens}[title={What you've built}]
One more word for telling what happened.
\end{grammarlens}

\subsection*{Your turn}
\begin{itemize}
\raggedright
  \item \emph{Say it:} \emph{wàijiè}
  \item \emph{Say it:} \emph{wàijiè}, once more
  \item \emph{Say it:} \emph{wàijiè dōu hěn guānxīn tā de bìng} — People everywhere are concerned about his illness.
  \item \emph{Recall:} say the Mandarin for a nation with its own government, then the Mandarin for the whole wide world, then say \emph{wàijiè} again
\end{itemize}

\begin{tcolorbox}[breakable,colback=teal!4,colframe=teal!35!black,title={Before you move on}]
What does \zh{外界} mean? (\textquotedblleft{}The outside world\textquotedblright{}.) Say it once more.
\end{tcolorbox}
\section[jiànlì]{jiànlì — to found; to set up}

\label{lesson:ZH-C326-jianli}



\begin{quote}
Before the new one: say the Mandarin for the whole wide world, then the Mandarin for the outside world.
\end{quote}

\subsection*{You'll want to know: jiànlì}
\emph{jiànlì} — \textquotedblleft{}to found; to set up\textquotedblright{}.

\begin{quote}
\emph{tāmen yìbǎi nián qián jiànlì le zhège xuéxiào} — They founded this school a hundred years ago.
\end{quote}

Its characters come later; for now it is a word for the ear.

\begin{cousinweb}
\emph{jiànlì} is built from \emph{jiàn} and \emph{lì}, both inherited from Middle Chinese, \emph{kjonH} and \emph{lip} (Guangyun, in Baxter's transcription), which glosses the first 'to set up' and 'to plant', and the second 'to stand' and 'to complete'.
\end{cousinweb}

\begin{grammarlens}[title={What you've built}]
One more word for telling what happened.
\end{grammarlens}

\subsection*{Your turn}
\begin{itemize}
\raggedright
  \item \emph{Say it:} \emph{jiànlì}
  \item \emph{Say it:} \emph{jiànlì}, once more
  \item \emph{Say it:} \emph{tāmen yìbǎi nián qián jiànlì le zhège xuéxiào} — They founded this school a hundred years ago.
  \item \emph{Recall:} say the Mandarin for the whole wide world, then the Mandarin for the outside world, then say \emph{jiànlì} again
\end{itemize}

\begin{tcolorbox}[breakable,colback=teal!4,colframe=teal!35!black,title={Before you move on}]
What does jiànlì mean? (\textquotedblleft{}To found; to set up\textquotedblright{}.) Say it once more.
\end{tcolorbox}
\section[\texorpdfstring{\zh{文明} (wénmíng)}{文明 (wénmíng)}]{\zh{文明} (wénmíng) — civilization}

\label{lesson:ZH-C326-wenming}



\begin{quote}
Before the new one: say the Mandarin for the outside world, then the Mandarin for to found something new.
\end{quote}

\subsection*{You'll want to know: \zh{文明}}
\textbf{\zh{文明}} — \emph{wénmíng} — \textquotedblleft{}civilization\textquotedblright{}.

\begin{quote}
\textbf{\zh{中国有很古老的文明}} \emph{Zhōngguó yǒu hěn gǔlǎo de wénmíng} — China has a very ancient civilization.
\end{quote}

\begin{cousinweb}
\textbf{\zh{文明}} (\emph{wénmíng}) is built from \zh{文} 'writing' and \zh{明} 'bright', both inherited from Middle Chinese, \emph{mjun} and \emph{mjaeng} (Guangyun, in Baxter's transcription), which glosses the first 'writing' and 'beautiful' and the second 'light'. The sense 'civilization' is not given there.
\end{cousinweb}

\begin{grammarlens}[title={What you've built}]
One more word for telling what happened.
\end{grammarlens}

\subsection*{Your turn}
\begin{itemize}
\raggedright
  \item \emph{Say it:} \emph{wénmíng}
  \item \emph{Say it:} \emph{wénmíng}, once more
  \item \emph{Say it:} \emph{Zhōngguó yǒu hěn gǔlǎo de wénmíng} — China has a very ancient civilization.
  \item \emph{Recall:} say the Mandarin for the outside world, then the Mandarin for to found something new, then say \emph{wénmíng} again
\end{itemize}

\begin{tcolorbox}[breakable,colback=teal!4,colframe=teal!35!black,title={Before you move on}]
What does \zh{文明} mean? (\textquotedblleft{}Civilization\textquotedblright{}.) Say it once more.
\end{tcolorbox}
